\chapter{The forecast model}
\label{ch:model}

\section{Task}
As in \citet{bian2025}, the model looks at the \cotwo{} map of the last 6~minutes (12 maps, 30~s apart) and
forecasts the \cotwo{} map of the next 3~minutes (6 maps). The maps are taken on the horizontal plane at
1.6~m, the height of the heads of standing people in the paper. On the 8~cm grid this plane is the cell layer
at 1.58~m with 21,956 points. Besides the \cotwo{} history, the model receives the speeds of the 8 windows and
the number of people.

\section{Network}
The network follows the GNOT idea \citep{hao2023gnot} (\tref{tab:model}). Every point where a forecast is wanted
(query) is described by its position, encoded with Fourier features, and its own \cotwo{} history. A fixed set
of 900 randomly chosen points of the plane carries the \cotwo{} history of the whole room, and the window speeds
and the number of people enter as additional input. In three attention blocks the query points collect
information from these inputs. For every query point the network outputs, for each of the 6 forecast times,
the change from the last known \cotwo{} value and the logarithm of its variance, so the model also gives an
uncertainty.

\begin{table}[h]
  \centering
  \caption{The forecast model.}
  \label{tab:model}
  \begin{tabular}{>{\raggedright}p{4cm}p{9.7cm}}
    \toprule
    Item & Setting\\
    \midrule
    Input & \cotwo{} of the last 6~min on the 1.6~m plane (12 maps), 8 window speeds, number of people\\
    Output & \cotwo{} change for the next 3~min (6 maps) with variance\\
    Network & cross-attention blocks (GNOT type), 3 blocks, width 128, 4 heads, about 480,000 parameters\\
    Loss & Gaussian negative log-likelihood, as in the paper\\
    Optimiser & AdamW, learning rate $5\times10^{-4}$ with one-cycle schedule, 20,000 steps,
               8 samples $\times$ 1,024 query points per step\\
    Ensemble & 5 networks with different random seeds; the forecast is their mean\\
    Hardware, time & NVIDIA RTX A2000 (12~GB), about 1.5~h for the 5 networks\\
    \bottomrule
  \end{tabular}
\end{table}

\section{Training data}
Each sample is a window of 9~minutes from one 30-minute run: 6~minutes of history and 3~minutes to forecast.
With a new sample every 30~s this gives 44 samples per run. Two kinds of runs are used:
\begin{itemize}
  \item \emph{Direct runs}: the room starts with fresh air, and the \cotwo{} of the people builds up.
  \item \emph{Transition runs} (paper-like): the room starts with the \cotwo{} of another setting
        (\cref{ch:dataset}); the numbers of people before and after the change are drawn between 10 and 80.
\end{itemize}

\section{Error measure}
The error of one sample is the relative $\ell_2$ error of the forecast over the plane and the 6 forecast times
(eq.~12 in \citet{bian2025}), and the reported error is the mean over all test samples. I report it in two
ways. The strict version uses the \cotwo{} above the outdoor level. The paper's version uses the total \cotwo{}
including the 400~ppm of the outdoor air; since these 400~ppm are the same in the forecast and in the data, they
make the error look much smaller.

As a baseline I use the simplest possible forecast: the \cotwo{} in 3~minutes is the same as now
(persistence). A useful model has to be clearly better than this. Because our open windows exchange the air of
the room much faster than the ceiling vents in the paper, the \cotwo{} changes differently than in their room,
and the comparison with the baseline on the same data is the fairest measure.
